% =====================================================================
% DRAFT — additional paragraphs for ch5 §"Relation to prior art, validation, and limits"
% (\subsection{Relation to prior art, validation, and limits}, label sec:fem_validation)
% Thesis body (30_Masterarbeit) is FROZEN — this is a draft to be slotted in on approval.
%
% Focus: the honest DATA-LAYER SEPARATION (patient composition is 0D; depth shape is in-vitro)
% and the bridge-assumption validation (load is the lever, borrowed shape is a thin band).
%
% New assets referenced (already built):
%   figures/F1b_data_bridge.pdf                  -> \label{fig:fem_databridge}
%   figures/F1c_patient_bridge_validation.pdf    -> \label{fig:fem_patientbridge}
% Generators: masterarbeit_ansys_fem/gen_bridge_fig.py ; extensions/c4_patient_bridge_validation.py
% Existing labels reused: sec:fem_residual, sec:fem_validation, sec:fem_tooth
% =====================================================================

\paragraph{Two real datasets, two different layers.} The framework is driven by real
data at two layers that must not be conflated. The patient cohort
(Dieckow~\textit{et al.}~peri-implant 16S~\cite{dieckow2024}) is real but
\emph{composition-only}: it resolves \emph{how much} of each guild is present, with no
depth axis, and it drives the ecological model of Chapters~\ref{ch:heine}--\ref{ch:dieckow},
not the finite-element model. The \emph{depth structure} the mechanics needs --- the
profile $\phi_\mathrm{Pg}(z)$ that sets the eigenstrain --- comes instead from the
in-vitro HOBIC confocal imaging~\cite{Heine2025PeriImplant}, where the
\textit{P.~gingivalis} channel is recovered by the colocalisation decode of
Section~\ref{sec:fem_growth} ($\mathrm{Fn}=\mathrm{blue}\cap\mathrm{red}$,
$\mathrm{Pg}=\mathrm{red}-\mathrm{Fn}$), not a one-to-one channel map. The two layers and
their join are shown in Figure~\ref{fig:fem_databridge}: a measured patient composition
(0D) supplies a Pg \emph{load}, the in-vitro confocal biofilm supplies a depth
\emph{shape}, and their product is the assumed patient profile
$\phi_\mathrm{Pg}(z)=\text{load}\times\text{shape}$. This is an explicit assumption, not a
per-patient measurement: no patient in the cohort was imaged in depth, so the depth shape
is \emph{borrowed} from the in-vitro model.

\begin{figure}[htbp]
  \centering
  \includegraphics[width=0.98\textwidth]{F1b_data_bridge.pdf}
  \caption[Data-provenance bridge: composition to depth to stress]{The two real datasets
  enter at different layers. \textbf{(A)}~real patient 16S (Dieckow, ten peri-implant
  patients) is composition-only --- a Pg-guild abundance with no depth.
  \textbf{(B)}~the depth \emph{shape} $\phi_\mathrm{Pg}(z)$ is measured on the in-vitro
  HOBIC biofilm (corrected confocal decode) and scaled by the patient load.
  \textbf{(C)}~the material model turns the assumed profile into a residual-stress column
  $S_{11}(z)/\mu$. The join in (B), $\text{load}\times\text{shape}$, is the bridge
  assumption examined below.}
  \label{fig:fem_databridge}
\end{figure}

\paragraph{The borrowed shape is a weak lever; the measured load is the strong one.}
Because the depth shape is borrowed rather than measured, the decisive question is how
much it matters --- and projecting the real cohort through the verified mechanics answers
it. Using the boundary-layer detachment driver
$J=\langle e^{-z/z_0}\,\phi_\mathrm{Pg}(z)\rangle_z$ ($z_0=0.25$), the verified reduced
form of the interface stress (only near-substratum growth loads the interface;
Section~\ref{sec:fem_residual}, where $S_{11}$ is depth-dominated rather than a local-$\phi$
law), each of the ten Dieckow patients maps to a detachment risk
(Figure~\ref{fig:fem_patientbridge}). The inter-patient spread is large --- a coefficient
of variation of $0.86$ across $J\in[0,\,0.028]$ --- and it is set entirely by the
\emph{measured} Pg load. The unmeasured depth shape, by contrast, is a thin band:
substituting the \emph{entire} measured commensal (CH) shape for the dysbiotic (DH) shape
at fixed load changes $J$ by only $8\%$, so that the inter-patient spread is
\emph{$34\times$ wider than the shape-uncertainty band}. The borrowing on which the bridge
rests is therefore low-risk: per-patient mechanical risk is governed by the quantity the
clinic can actually measure (the 0D Pg load), not by the per-patient depth structure the
model cannot see.

\begin{figure}[htbp]
  \centering
  \includegraphics[width=0.98\textwidth]{F1c_patient_bridge_validation.pdf}
  \caption[Patient-bridge validation: load is the lever, not shape]{Projecting the real
  Dieckow cohort through the verified detachment kernel
  $J=\langle e^{-z/z_0}\phi_\mathrm{Pg}(z)\rangle$. \textbf{(A)}~per-patient risk spans a
  factor that the measured Pg load alone explains (CV~$=0.86$); the error bars are the
  shape-uncertainty band obtained by swapping the measured CH and DH depth shapes.
  \textbf{(B)}~the bridge law $J=\text{load}\times\text{shape}$ is linear; the grey wedge
  is the $\pm8\%$ shape band. The inter-patient spread is $34\times$ the shape band, so the
  borrowed in-vitro shape is a weak lever relative to the measured load.}
  \label{fig:fem_patientbridge}
\end{figure}

\paragraph{What this does and does not license.} Two boundaries keep the claim honest.
The low shape-leverage is established for the boundary-layer \emph{driver} $J$, which ranks
detachment risk; it does \emph{not} extend to the absolute stress \emph{scale}, which
remains parametric in the biofilm modulus for the reasons given above (the
shape-calibrated, scale-parametric reading of Section~\ref{sec:fem_validation}). And the
result says the per-patient depth \emph{redistribution} is a small correction within a
unit-mean envelope --- not that depth is irrelevant: the depth \emph{kernel} itself, which
concentrates the interface stress near the substratum, is what makes the mechanism
composition-resolved in the first place (Section~\ref{sec:fem_residual}), and the
geometry-portable real-tooth run of Section~\ref{sec:fem_tooth} likewise still relies on
the in-vitro depth shape. The validation that would remove the borrowing altogether is the
same experiment identified above, now sharpened: \emph{depth-resolved imaging of a patient
biofilm} would replace the assumed shape with a measured one, and per-patient
\emph{shear-flow detachment onset versus Pg load} would test the linear $J$--load law that
the cohort projection predicts.
